% A143: actual elliptic quotient, exact rank, and profinite lifting exclusion.
\subsection{The rank-two elliptic quotient and its excluded profinite classes}

\begin{statement}{A.143}{Proved}
{An explicit positive-rank factor with a finite-primary lifting obstruction}
For the proper genus-three curve \(\overline D\) of A.139 put
\[
 \begin{gathered}
 k=573,\qquad n=1146=2k,\qquad
 q_0=\frac{325596479663}{72962},\qquad a=kq_0^2,\\
 E:\ y^2=x^3-n^2x,\qquad F:\ v^2=u^3+k^2u.
 \end{gathered}
\tag{A.143.1}
\]
There are actual morphisms over \(\Q\)
\[
              \overline D\xrightarrow{\ q\ }E
                         \xrightarrow{\ \phi\ }F
\tag{A.143.2}
\]
of degrees two and two.  Their composite is the quotient by the
rational order-four automorphism
\[
 r(T,W)=\left(\frac{T-1}{2T-1},\,\frac{W}{2T-1}\right).
\tag{A.143.3}
\]
Both elliptic curves have rational Mordell--Weil rank exactly two,
and their central \(L\)-values are zero.
Two independent infinite-order points on \(E\) are
\[
                   P=(2028,75348),\qquad Q=(5202,365976).
\tag{A.143.4}
\]
Their images on \(F\) are
\[
 \phi(P)=\left(\frac{233289}{676},\frac{218400525}{17576}\right),
 \qquad
 \phi(Q)=\left(\frac{357604}{289},\frac{235662830}{4913}\right).
\tag{A.143.5}
\]
Moreover \(\phi:E(\Q)\to F(\Q)\) is onto.

Let \(\eta\in\widehat{F(\Q)}\) be a non-torsion element whose
primary components are zero outside a finite set of coefficient
primes.  Its proper elliptic Kummer section does not have a full
Selmer lift to \(D^\circ\) through (A.143.2).
In particular, for every non-torsion \(R\in F(\Q)\) and every
nonzero \(z\in\widehat\Z\) supported at finitely many coefficient
primes, the class \(zR\) has no such lift.  The support may include
the coefficient prime two.

This is a positive-rank conclusion on the displayed curve and map.
No value of the Tate--Shafarevich group, exact positive analytic
rank, or classification of the infinite-primary-support classes
is asserted.
\end{statement}

\begin{proof}
\emph{The quotient morphism and its isogeny factorization.}
Set \(s=2T-1\), so that
\[
                         W^4=\frac{s^3-s}{8a}.
\tag{A.143.6}
\]
The automorphism (A.143.3) becomes
\[
                         (s,W)\longmapsto(-1/s,W/s).
\tag{A.143.7}
\]
Substitution proves it preserves (A.143.6); its square is
\((s,W)\mapsto(s,-W)\), so it has order four.
The invariant functions
\[
                    \xi=s-s^{-1},\qquad
                    z=W^2(s^2+1)/s^2
\tag{A.143.8}
\]
satisfy
\[
                         z^2=\frac{\xi(\xi^2+4)}{8a}.
\tag{A.143.9}
\]
Conversely they recover the function field by
\[
                     s^2-\xi s-1=0,\qquad
                     W^2=\frac{zs^2}{s^2+1}.
\tag{A.143.10}
\]
These equations give degree at most four; the four distinct powers
of (A.143.7) fix the two functions, so the degree is exactly four.
The extension is the actual invariant-field extension.
All rational maps here extend uniquely to the smooth proper curves.
After the exact scaling \(a=kq_0^2\), the quotient coordinates are
\[
                    u=\frac{k}{2}(s-s^{-1}),\qquad
                    v=k^2q_0 W^2(s^2+1)/s^2,
\tag{A.143.11}
\]
and they satisfy the equation of \(F\).

More precisely, the intermediate degree-two map in (A.143.2) is
\[
                         x=ns,\qquad
                         y=8k^2q_0W^2 .
\tag{A.143.12}
\]
Its equation is verified directly using (A.143.6).
Its involution is \(W\mapsto-W\), and \(T\) and \(W^2\) are
recovered from \(x,y\), proving degree two.
The two mutually dual isogenies used below are
\[
 \begin{aligned}
 \phi:E\longrightarrow F,\qquad
  (x,y)&\longmapsto
    \left(\frac{x-n^2/x}{4},
           \frac{y(1+n^2/x^2)}8\right),\\
 \psi:F\longrightarrow E,\qquad
  (u,v)&\longmapsto
    \left(u+\frac{k^2}{u},
           v(1-k^2/u^2)\right).
 \end{aligned}
\tag{A.143.13}
\]
They extend over \(x=0,u=0\) and infinity, send origins to
origins, and have kernels generated by \((0,0)\).
Their curve equations follow by substitution.
A nonconstant morphism of elliptic curves taking origin to origin
is a group homomorphism: under
\(P\mapsto[P-O]\), its pushforward on degree-zero divisor classes
sends this class to \([f(P)-O]\), and principal divisors push
forward by the function-field norm.
Thus this also proves the homomorphism assertion for these
particular morphisms.
Their degrees are two by the quadratic inverse equations implicit
in (A.143.13).

Substituting (A.143.12) in \(\phi\) gives exactly (A.143.11).
The automorphism of \(E\) induced by \(r\) in (A.143.7) is
\[
                         (x,y)\longmapsto
                  (-n^2/x,\ n^2y/x^2),
\tag{A.143.14}
\]
which is the chord-law translation by \((0,0)\).
This proves the factorization, including its quotient maps.
Both compositions in (A.143.13) are multiplication by two.
For example the \(E\)-composition gives
\[
 \begin{gathered}
 x_2=\frac{(x^2+n^2)^2}{4y^2},\qquad
 y_2=\frac{3x^2-n^2}{2y}(x-x_2)-y ,
 \end{gathered}
\tag{A.143.15}
\]
which is the literal duplication formula.
The corresponding substitution on \(F\) gives its duplication
formula as well.

For later use the fibre over the two rational boundary images is
exactly
\[
 \begin{aligned}
 (\phi q)^{-1}(O_F)
      &=\{T=\tfrac12,\ T=\infty\},\\
 (\phi q)^{-1}((0,0)_F)
      &=\{T=0,\ T=1\}.
 \end{aligned}
\tag{A.143.16}
\]
Indeed \(\phi^{-1}(O_F)=\{O_E,(0,0)_E\}\), while
\(\phi^{-1}((0,0)_F)=\{(n,0),(-n,0)\}\).
Equation (A.143.12) sends the four boundary points to exactly
these four two-torsion points.  The degree-two map \(q\) is
ramified with index two at each: \(W\mapsto-W\) fixes precisely
the three finite \(W=0\) points and the point at infinity.
Since \(\phi\) is etale in characteristic zero, every point in
(A.143.16) has ramification index two.
The four degree counts exhaust both fibres.
Consequently
\[
             D^\circ\longrightarrow
                 F\setminus\{O_F,(0,0)_F\}
\tag{A.143.17}
\]
is the actual degree-four finite etale map.

\emph{An elementary exact two-isogeny descent.}
For a curve \(Y^2=X(X^2+b)\) define
\[
 \delta_b(O)=1,\qquad \delta_b((0,0))=[b],\qquad
                    \delta_b((X,Y))=[X]\quad(X\ne0)
\tag{A.143.18}
\]
in \(\Q^\times/(\Q^\times)^2\).
This is a homomorphism with the indicated formulas.
For a line \(Y=mX+c_0\), the product of the three nonzero
intersection \(X\)-coordinates is \(c_0^2\); the chord rule
therefore proves the assertion, including tangent lines by
multiplicity.  Addition of \((0,0)\) has new coordinate \(b/X\),
which gives its prescribed factor \([b]\).
Vertical lines and the origin give the remaining cases.

Write \(\delta_E=\delta_{-n^2}\) and
\(\delta_F=\delta_{k^2}\).
Their kernels are exactly
\[
                     \ker\delta_E=\psi(F(\Q)),\qquad
                     \ker\delta_F=\phi(E(\Q)).
\tag{A.143.19}
\]
One inclusion follows from
\(u+k^2/u=(v/u)^2\) on \(F\) and
\((x-n^2/x)/4=(y/(2x))^2\) on \(E\).
For the converse on \(E\), write \(x=r_0^2\ne0\) and set
\[
             u=\frac{x+y/r_0}{2},\qquad v=r_0u .
\tag{A.143.20}
\]
Then \(u\ne0\), \(u^2-xu+k^2=0\), the point lies on \(F\),
and \(\psi(u,v)=(x,y)\).
For the converse on \(F\), write \(u=r_0^2\ne0\) and set
\[
                         x=2u+2v/r_0,\qquad y=2r_0x .
\tag{A.143.21}
\]
Then \(x^2-4ux-n^2=0\), this point lies on \(E\), and
\(\phi(x,y)=(u,v)\).
The origin is always in the image.
The remaining point \((0,0)_F\) has square class one and is
the image of \((\pm n,0)_E\);
\((0,0)_E\) has class \(-1\) and is not in the kernel.
These verify every exceptional case in (A.143.19).

At a prime not dividing \(b\), a nonzero rational \(X\) on
\(Y^2=X(X^2+b)\) has even valuation.
For positive valuation the right side has valuation \(v_p(X)\);
for negative valuation it has valuation \(3v_p(X)\).
Both must be even.  Thus the image of \(\delta_E\) is contained
in the sixteen squareclasses generated by \(-1,2,3,191\).
All sixteen occur.  The points
\[
 \begin{array}{c|cccc}
 \text{point}&(0,0)&(1146,0)&P&Q\\ \hline
 \delta_E&-1&1146&3&2
 \end{array}
\tag{A.143.22}
\]
give four independent squareclasses: \(1146=2\cdot3\cdot191\).
The equations of \(P,Q\) are verified by integer substitution.
The exact constants also explain their origins.
On the intermediate equation \(Z^2=F(T)/a\), the old centre
\(T=R\) has \(Z=1/155564491\), giving \(P\) under
(A.143.12) with \(W^2=Z\).
At \(T=c\), choose \(Z=-i\rho\).
Its \(E\)-image is \((-5202,-365976i)\).
The actual automorphism \((x,y)\mapsto(-x,iy)\) gives \(Q\).

For \(F\), every nonzero real point has \(u>0\).
The valuation argument therefore leaves only the positive classes
\[
                                  1,\ 3,\ 191,\ 573 .
\tag{A.143.23}
\]
A point of class \(d_0\) among the last three would give, after
writing \(u=d_0(U/V)^2\) in lowest terms, the primitive equation
\[
                 N^2=d_0U^4+\frac{k^2}{d_0}V^4,
                      \qquad \gcd(U,V)=1 .
\tag{A.143.24}
\]
Here \(N=vV^3/(d_0U)\) is rational with integral square, hence
integral; \(u\ne0\).
For \(d_0=3\) the right side is
\(3(U^4+191^2V^4)\), and for \(d_0=573\) it is
\(573(U^4+V^4)\).
Unless both \(U,V\) are divisible by three, the parenthesis is
nonzero modulo three: fourth powers are zero or one there.
The valuation is exactly one, impossible for \(N^2\).
Primitivity excludes the remaining alternative.
For \(d_0=191\), the equation becomes
\(N^2=191(U^4+9V^4)\).
Modulo \(191\), a vanishing parenthesis with one nonzero
coordinate would make \(-9\) a fourth power.  It is not even
a square, since \(191\equiv3\pmod4\).
Again both coordinates would have to be divisible by \(191\),
or the right side would have odd valuation.  Both alternatives
are impossible.
We have proved
\[
                 |\operatorname{im}\delta_E|=16,\qquad
                 \operatorname{im}\delta_F=\{1\}.
\tag{A.143.25}
\]

It follows from (A.143.19) that \(F(\Q)=\phi(E(\Q))\).
Using \(\psi\phi=[2]\) gives
\[
                        |E(\Q)/2E(\Q)|=16.
\tag{A.143.26}
\]
Apply the precise height form of the finite-generation argument
to the sixteen-element quotient just proved.
Choose one rational representative of each of its sixteen
cosets, and choose \(C>0\) larger than their canonical heights.
The set of rational points of canonical height at most \(C\)
is finite by \cite[IV, Theorem 4.7(a)]{AgentMilneEC}.
It contains these coset representatives and therefore generates
\(E(\Q)\) by \cite[IV, Proposition 4.13]{AgentMilneEC}.
This proves the needed finite generation from (A.143.26).
Its rational two-torsion is exactly
\(\{O,(0,0),(\pm1146,0)\}\), of order four.
For a finitely generated abelian group, the order of its quotient
by twice the group is \(2^{\rm rank}\) times the order of its
two-torsion.  Hence the rank is two.
The isogenies induce inverse isomorphisms up to multiplication
by two after tensoring with \(\Q\), so \(F\) has rank two too.
The images of the rational two-torsion under \(\delta_E\)
already give the four classes \(\{1,-1,1146,-1146\}\).
The classes of \(P,Q\) are independent modulo these.
Since the torsion subgroup modulo twice itself has order four,
this proves that \(P,Q\) are independent modulo torsion.
Their images in (A.143.5) are obtained by (A.143.13).
No torsion-integrality theorem or rank conjecture is needed.

\emph{The exact central \(L\)-value calculation.}
Use \(\mathfrak q\) as a formal series variable.
Tunnell's Theorem 1 gives the signed theta expansion
\[
 g=\sum_{m,j\in\Z}(-1)^j
                 \mathfrak q^{(4m+1)^2+8j^2},
      \qquad
 g\theta_4=\sum_{\nu\ge1}b(\nu)\mathfrak q^\nu,
      \quad \theta_4=\sum_{\ell\in\Z}\mathfrak q^{4\ell^2}.
\tag{A.143.27}
\]
The opened original paper is
\cite[Theorem 1, p.~326; Theorem 3, pp.~328--329]{LegTunnell1983}.
Exhaustive integer enumeration at \(\nu=573\) gives sixteen
terms with sign \(+1\) and sixteen with sign \(-1\).
Thus \(b(573)=0\).
Equivalently, the full ternary counts are
\[
 \begin{aligned}
 \#\{(X,Y,Z)\in\Z^3:573=4X^2+Y^2+8Z^2\}&=64,\\
 \#\{(X,Y,Z)\in\Z^3:573=4X^2+Y^2+32Z^2\}&=32.
 \end{aligned}
\tag{A.143.28}
\]
For odd \(Y\), exactly one of \(Y,-Y\) is \(1\) modulo four.
The signed sum in (A.143.27) is consequently the second count
minus one half of the first, explaining their literal comparison.
The enumeration ranges are bounded by the square roots of the
positive summands, so they are exhaustive.
The retained code separately verifies the same coefficient in
\(\mathfrak q\prod_{m\ge1}(1-\mathfrak q^{8m})
(1-\mathfrak q^{16m})\theta_4\).

Since \(573=3\cdot191\) is odd and squarefree, Tunnell's exact
formula applies and gives
\[
 L(E,1)=\frac{b(573)^2\beta}{2\sqrt{1146}}=0,\qquad
       \beta=\int_1^\infty\frac{dx}{\sqrt{x^3-x}}>0 .
\tag{A.143.29}
\]
The isogenies identify the rational Tate representations, since
their composites are multiplication by two.  Their local Euler
factors, including the factors on inertia invariants, therefore
agree.  Hence \(L(F,1)=L(E,1)=0\).
This proves the central-value assertion without a converse to
Tunnell's theorem and without an assertion about the exact order
of vanishing.

\emph{Finite-primary completion classes do not lift to the punctured curve.}
Let \(S\) be a finite set of coefficient primes containing the
support of \(\eta\), and let \(e_S\in\widehat\Z\) be its
primary idempotent.
The surjection \(E(\Q)\to F(\Q)\) induces a surjection on
profinite completions: the compact image of
\(\widehat{E(\Q)}\) contains the dense image \(F(\Q)\), and is
therefore all of \(\widehat{F(\Q)}\).
Choose a preimage and multiply it by \(e_S\).
This gives
\[
              \theta=e_S\theta\in\widehat{E(\Q)},\qquad
                                \widehat\phi(\theta)=\eta.
\tag{A.143.30}
\]
Localization is a continuous homomorphism from these completions
to the compact groups \(E(\Q_p)\) and \(F(\Q_p)\);
it commutes with all primary idempotents.
Their full local Kummer maps are injective: each group has an
open pro-\(p\) subgroup and finite quotient, so its elements are
separated by the quotients by \(m\)-multiples.

Take good physical primes outside \(S\) with
\[
                   p\equiv3\pmod8,\qquad
                   p\equiv1\pmod\ell
                         \quad(\ell\in S,\ \ell\ \text{odd}).
\tag{A.143.31}
\]
These compatible conditions specify a nonempty Frobenius class
in a cyclotomic extension, hence a set of positive Dirichlet
density.  Excluding \(2,3,191\) removes only finitely many primes.
For either elliptic equation the defining cubic is odd.
Because \(-1\) is a nonsquare modulo \(p\), pairing \(x\) with
\(-x\) cancels its quadratic-character sum.  Therefore
\[
                    \#E(\F_p)=\#F(\F_p)=p+1 .
\tag{A.143.32}
\]
Every odd \(\ell\in S\) is prime to \(p+1\).
The corresponding primary local component of \(\theta\) is
zero, since the reduction kernel is pro-\(p\).
Moreover \(v_2(p+1)=2\), while all four points of \(E[2]\)
are rational and reduce distinctly.
The full two-primary subgroup of \(E(\Q_p)\) is thus exactly
\(E[2]\): reduction identifies its prime-to-\(p\) torsion
with that in the finite group.
Consequently
\[
                     \eta_p\in\phi(E[2])
                          =\{O_F,(0,0)_F\}
\tag{A.143.33}
\]
at every prime in (A.143.31), whether or not \(2\in S\).

Suppose that a full Selmer section of \(D^\circ\) has this
proper elliptic Kummer class.
At any prime in (A.143.31), local Kummer injectivity and
(A.143.16) show that its local centre is a boundary point:
an ordinary point of \(D^\circ\) cannot map into the displayed
two-point set.  Push this full section through the actual
finite etale map \(D^\circ\to U\) of A.139.
It is locally cuspidal at this positive-density set.
The density-one integrality theorem AP-D.25 in its Selmer range,
used with the exact same scope in A.23, says that a globally
noncuspidal image has ordinary integral local points on a
density-one set \cite[D.25]{LegAP2026}.
The two assertions are incompatible.  The image on \(U\) is
therefore globally cuspidal.

For the transfer back, identify \(E_D\) with the open subgroup
associated to the actual degree-eight cover.
The intersection of a base cusp decomposition group with this
subgroup is the decomposition group of an extension of that
cusp valuation in \(\overline D\).
After conjugation the section image lies in this intersection.
Its surjection onto \(G_\Q\) makes the corresponding boundary
point rational.  Thus it is globally cuspidal on \(D^\circ\).
Compactification kills its inertia, and (A.143.16) sends its
proper elliptic class to the rational torsion class of
\(O_F\) or \((0,0)_F\).

The global Kummer map on \(\widehat{F(\Q)}\) is injective:
the finite maps \(F(\Q)/mF(\Q)\to H^1(\Q,F[m])\) are injective,
and their compatible inverse-limit map remains injective.
Thus the non-torsion \(\eta\) cannot equal either torsion class.
This contradiction proves the lifting exclusion.
These compact classes themselves define continuous proper
elliptic sections, using
\(\pi_1(F)=T(F)\rtimes G_\Q\) from the origin and their Kummer
cocycles.  Explicitly, choose finitely many rational
Mordell--Weil generators and compatible division points for them;
their continuous Tate-module cocycles and their
\(\widehat\Z\)-linear combinations give the completion classes.
For the explicit scalar family, multiplication of a
Kummer cocycle of \(R\) by \(z\) is the literal cocycle giving
the proper section of \(zR\).
Nonzero \(z\) of finite primary support preserves a nonzero
free primary component of a non-torsion \(R\), so \(zR\) is
non-torsion.  This proves the stated scalar consequence.
Only the Selmer hypothesis, not an additional local
classification theorem, was used.
\end{proof}

The code \path{code/harmonic_elliptic_rank_certificate.py},
its retained output and \path{code/harmonic_elliptic_rank_data.json}
check the curve morphisms, both isogenies and both duplication
identities, the exact points and their images, all \(16\)
squareclasses, and every primitive residue pair in the three
quartic exclusions.  They retain all \(32\) signed theta terms,
all \(96\) ternary representations, and the independent product
coefficient.  A separate finite-field check computes the full
two-primary projections and their isogeny images at good primes
congruent to three modulo eight.
The exact doubling value
\(x(2P)=22718025625/17522596\) is also recorded; the rank proof
uses the explicit descent rather than a torsion-integrality test.
